
%%% Обязательные пакеты
%% Beamer
\usepackage{beamerthemesplit}
\usetheme{SPbGU}
\beamertemplatenavigationsymbolsempty
\usepackage{appendixnumberbeamer}

%% Локализация
\usepackage{fontspec}
\setmainfont{PT Serif}
\setsansfont{Helvetica}
\setmonofont{Menlo}

\usepackage{polyglossia}
\setdefaultlanguage{russian}
\setotherlanguage{english}
\newfontfamily\russianfont{PT Serif}
\newfontfamily\russianfontsf{Helvetica}
\newfontfamily\russianfonttt{Menlo}
\usepackage[autostyle]{csquotes}

%% Графика
\usepackage{wrapfig}
\usepackage{pdfpages}
\usepackage{pgfplots}
\pgfplotsset{compat=1.18}

%% Математика
\usepackage{amsmath, amsfonts, amssymb, amsthm, mathtools}

% Математические окружения с русским названием
\newtheorem{rutheorem}{Теорема}
\newtheorem{ruproof}{Доказательство}
\newtheorem{rudefinition}{Определение}
\newtheorem{rulemma}{Лемма}


%%% Дополнительные пакеты
\usepackage{tikz}
\usetikzlibrary{decorations.pathreplacing,calc,shapes,positioning,tikzmark,patterns}

\usepackage{multirow}

%% Пакеты для оформления алгоритмов на псевдокоде
\usepackage[noend]{algpseudocode}
\usepackage{algorithm}
\usepackage{algorithmicx}

\usepackage{fancyvrb}

%% Пакет для анимированных иллюстраций
\usepackage{animate}
